% No preamble. Prefix fwp-. Full positive Wilson measure, 17 October 2026.
\section{A full-measure Wilson Peierls bound}
\label{fwp-section}

We prove a rare-event estimate under the actual normalized positive
Wilson measure. It is independent of the Gaussian small-field analysis.
The event controlled here is large plaquette curvature; identifying it
with the complement of the actual retained-coordinate window is a
separate, unresolved interface.

Let \(G=SU(N)\), \(m_G=\dim G\), \(\tau=\operatorname{Tr}/N\), and
\[
 \Lambda=\prod_{i=1}^d\mathbb Z/n_i\mathbb Z,\qquad
 n_i=2m_i b,\quad m_i\ge1,\quad b\ge3.
\]
There are \(V=|\Lambda|\) sites, \(dV\) independent links and
\(\binom d2V\) elementary unoriented plaquettes. Normalized Haar measure
is used on every link. Set
\begin{equation}
 h(W)=1-\operatorname{Re}\tau(W)
       =\tfrac12\|W-I\|_\tau^2,\qquad
 A(U)=\sum_p h(U_p),\qquad
 \mu_g(dU)=Z_g^{-1}e^{-A(U)/g^2}\prod_e dU_e,
 \label{fwp-measure}
\end{equation}
where \(\|X\|_\tau^2=\tau(X^*X)\).
All side lengths are at least six, avoiding short-period graph
degeneracies.

Fix any translated grid of cubes of side \(b\).
Let \({\cal P}_j^{\rm int}\) be the plaquettes whose vertices are strictly
inside cell \(j\), and define
\begin{equation}
 {\cal B}_j(\delta)
 =\left\{\max_{p\in{\cal P}_j^{\rm int}}h(U_p)\ge\delta^2\right\}.
 \label{fwp-event}
\end{equation}
Strict interior means that every coordinate of every used link endpoint
lies strictly between the corresponding walls.

\begin{theorem}[Prescribed bad cells]
\label{fwp-main}
There are \(g_0>0\) and \(C_{G,d}<\infty\), depending only on the group
and dimension, such that for \(0<g\le g_0\) and every set \(S\) of
distinct cells,
\begin{equation}
 \mu_g\left(\bigcap_{j\in S}{\cal B}_j(\delta)\right)
 \le z_b(g,\delta)^{|S|},
 \quad
 z_b(g,\delta)=
 \min\left\{1,
 \exp\left[-\frac{\delta^2}{g^2}
     +b^d\{d m_G\log(1/g)+C_{G,d}\}\right]\right\}.
 \label{fwp-main-bound}
\end{equation}
Constants are uniform in volume and in the selected cells. The source
regulators \(n=2L^M\), \(L\ge3\) odd, admit the construction with
\(b=L^u\) for each fixed \(u\ge1\) and \(M\ge u\).
\end{theorem}

\subsection{Exact site-plane slabs and the normalized chessboard inequality}
\label{fwp-slab-section}

Every cell wall is an integer coordinate, or \emph{site}, plane.
Its reflection preserves the lattice. An elementary plaquette cannot
have vertices on both strict sides of such a plane. It either lies in
the plane or belongs wholly to one adjacent closed half-space, possibly
touching the plane. Conditional on tangential links in the two opposite
reflection planes of the torus, the two open half configurations have
reflected product Haar measures and reflected Wilson weights.
The reflection quadratic form is therefore an integral of an absolute
square against a nonnegative boundary weight. This proves the needed
site-plane reflection positivity directly.

To obtain the chessboard inequality for an arbitrary even number of
cells, cut first in one coordinate into \(p=2m\) slabs of width \(b\).
Let \({\cal H}\) be \(L^2\) of the product Haar measure on tangential
links in one boundary plane. Integer translation identifies successive
boundary spaces.

For a bounded slab observable \(f\), supported strictly inside the
slab in the cut direction, form the integral kernel \(K_f\) by integrating
all nonboundary slab links. Give it the full action of every plaquette
assigned to the slab, including plaquettes touching a boundary but not
lying in it, and half the action of each plaquette lying entirely in
either boundary. Multiply by \(f\).
The compact integration domain and bounded integrand make \(K_f\)
Hilbert--Schmidt.

Gluing the kernels by boundary Haar integration gives exactly
\begin{equation}
 \operatorname{Tr}(K_{f_0}\cdots K_{f_{p-1}})
       =Z_g\!\left[\prod_{j=0}^{p-1}f_j\right].
 \label{fwp-gluing}
\end{equation}
Every boundary link is integrated once; every boundary plaquette gets
two halves; every other link and plaquette is assigned once.
No gauge-volume factor occurs. Reflection exchanges the two boundaries,
preserves \(h(U_p)\), and inverts a chosen link coordinate when its
orientation reverses. Haar invariance gives
\begin{equation}
 K_f^*=K_{\theta\overline f}.
 \label{fwp-adjoint}
\end{equation}
For real indicators complex conjugation has no effect.

Schatten H\"older, with the even integer \(p\), gives
\begin{equation}
 \left|\operatorname{Tr}\prod_{j=0}^{p-1}K_{f_j}\right|
 \le\prod_{j=0}^{p-1}
 \left[\operatorname{Tr}\bigl((K_{f_j}^*K_{f_j})^{p/2}\bigr)
                                                  \right]^{1/p}.
 \label{fwp-holder}
\end{equation}
Each factor is the exact Wilson partition integral with \(p\) alternating
reflected copies of one slab. The torus period remains \(pb\).
This explains precisely why the even cell count is required.

Apply \eqref{fwp-holder} successively in all coordinate directions.
A slab source is the product of its constituent cell sources.
Orthogonal affine coordinate reflections commute on the unoriented
graph, with the corresponding orientation inversions on link variables.
The final factor associated with an original cell is therefore its
full reflected dissemination. Intersections of cut planes require no
extra normalization: every step is the exact one-axis gluing
\eqref{fwp-gluing}.

Write \(N_c=V/b^d\) for the number of cells. Let \(\Theta_j\) map the
reference cell to cell \(j\) by the alternating reflection pattern;
even counts make this pattern periodic. For a reference-cell observable
\(f\), write
\({\cal D}(f)=\prod_i\Theta_i f\).
Pull each original observable back to this reference cell.
The result of the iteration is
\begin{equation}
 \left|Z_g\left[\prod_j\Theta_j f_j\right]\right|
 \le\prod_j Z_g[{\cal D}(f_j)]^{1/N_c}.
 \label{fwp-chessboard}
\end{equation}
For nonnegative indicators, all terms are nonnegative. Dividing by the
\emph{full} partition function, and using \(f_j=1\) off a selected set,
gives
\begin{equation}
 \mu_g\left(\prod_{j\in S}\Theta_j f_j\right)
 \le\prod_{j\in S}
       \left(\frac{Z_g[{\cal D}(f_j)]}{Z_g}\right)^{1/N_c}.
 \label{fwp-normalized}
\end{equation}
In particular, the normalization is not an independently chosen
small-field partition function.

The geometry used here is simpler than a diagonal fold: all walls are
site planes and no alternating site/link scheme is needed.
For \(n=2L^M\), \(b=L^u\), the number of cells in every direction is
\(2L^{M-u}\), which is even despite \(L\) being odd.

\subsection{A lower bound for the full normalization}
\label{fwp-normalization-section}

There are \(\epsilon_0>0\) and \(c_G\in(0,1]\) such that
\begin{equation}
 \operatorname{Haar}\{U:\|U-I\|_\tau\le\epsilon\}
       \ge c_G\epsilon^{m_G},\qquad 0<\epsilon\le\epsilon_0.
 \label{fwp-haar-ball}
\end{equation}
Indeed the Haar density in a fixed local exponential chart is smooth
and strictly positive. This is a finite-dimensional group estimate.

If all links lie in that \(\epsilon\)-ball, a four-term product
telescoping and unitarity show
\[
 \|U_p-I\|_\tau\le4\epsilon,\qquad h(U_p)\le8\epsilon^2.
\]
For \(g\le g_0:=\min(\epsilon_0,1)\), restrict the full integral to
\(\|U_e-I\|_\tau\le g\) on every link. It follows that
\begin{equation}
 Z_g\ge(c_Gg^{m_G})^{dV}
             \exp\!\left[-8\binom d2V\right],
 \qquad
 -\frac1V\log Z_g
       \le d m_G\log(1/g)+C_{G,d},
 \label{fwp-z-lower}
\end{equation}
where one can take
\[
 C_{G,d}=8\binom d2+d\log(1/c_G).
\]
This is only a lower bound for the normalization. It does not claim
that the small-link box has large probability.

On the full dissemination of \({\cal B}_j(\delta)\), every cell
contains a distinct interior plaquette of energy at least \(\delta^2\).
Thus \(A(U)\ge N_c\delta^2\) there and
\begin{equation}
 Z_g[{\cal D}(\mathbf1_{\cal B})]
       \le e^{-N_c\delta^2/g^2}.
 \label{fwp-disseminated-energy}
\end{equation}
The remaining product Haar integral is at most one.
Combining \eqref{fwp-z-lower} and \eqref{fwp-disseminated-energy}
bounds each factor in \eqref{fwp-normalized} by the exponential in
\eqref{fwp-main-bound}; it is also at most one.
This proves Theorem~\ref{fwp-main}.

More generally, a cell event forcing its strictly interior action to be
at least \(e_j\) has the factor
\begin{equation}
 \min\left\{1,\exp\left[-e_j/g^2
              +b^d\{d m_G\log(1/g)+C_{G,d}\}\right]\right\}.
 \label{fwp-energy-general}
\end{equation}
A threshold \(\|U_p-I\|_\tau\ge\delta\) uses \(e_j=\delta^2/2\).
At fixed \(b\), the threshold \(\delta(g)=gR(g)\) is rare whenever
\(R(g)^2/\log(1/g)\to\infty\), for example
\(R(g)=(\log(1/g))^a\) with \(a>1/2\).
The explicit \(b^d\) cost precludes claiming the same constants when
the cell size grows with the cutoff.

\subsection{Prescribed plaquettes and Peierls clusters}
\label{fwp-color-section}

Let \(S\) be any prescribed set of distinct elementary plaquettes.
Color \((x;\mu,\nu)\), \(\mu<\nu\), by \((x\bmod b,\mu,\nu)\).
There are
\begin{equation}
 Q_b=b^d\binom d2
 \label{fwp-color-count}
\end{equation}
colors. Choose deterministically a color occurring at least
\(|S|/Q_b\) times, with anchor residue \(a\), and shift the cell grid
by \(a-(1,\ldots,1)\). Every selected anchor has relative coordinates
\((1,\ldots,1)\); its plaquette vertices have coordinates \(1\) or \(2\)
and hence are strictly interior for \(b\ge3\).
Distinct anchors of the fixed color lie in distinct cells.
The fixed orientation prevents two selected plaquettes at one anchor.
The all-bad event on \(S\) implies the bad-cell event on this subset.
Therefore
\begin{equation}
 \mu_g\{h(U_p)\ge\delta^2\text{ for all }p\in S\}
       \le z_b(g,\delta)^{|S|/Q_b}.
 \label{fwp-arbitrary-plaquettes}
\end{equation}
This is a finite deterministic coloring argument, not an independence
assumption. The chessboard proof holds for each of the translated grids.

For the nearest-neighbor cell graph, of degree at most \(2d\), the
number of connected sets of size \(m\) containing a fixed cell is at
most \((2d)^{2m}\): choose a canonical spanning tree and its depth-first
traversal. The traversal has length \(2(m-1)\) and determines the visited
set. A cluster of size at least \(m\) contains such a connected subset
of size exactly \(m\). Consequently
\begin{equation}
 \mu_g\{|\text{bad cluster at }j|\ge m\}
       \le[(2d)^2z_b(g,\delta)]^m.
 \label{fwp-cell-cluster}
\end{equation}
For any fixed finite-degree plaquette adjacency graph the same argument
and \eqref{fwp-arbitrary-plaquettes} apply. For adjacency by shared
vertices, \(\Delta_d=16\binom d2\) is a safe degree bound, giving
\begin{equation}
 \mu_g\{|\text{bad plaquette cluster at }p|\ge m\}
       \le[\Delta_d^2z_b(g,\delta)^{1/Q_b}]^m.
 \label{fwp-plaquette-cluster}
\end{equation}
Periodic identifications do not increase these counting bounds.
These are intersection and cluster estimates, not a claim of full
stochastic domination by a product measure.

The union estimate for any bad cell is only \(N_c z_b\).
Thus the theorem does not assert that an arbitrarily large torus lies
entirely in one small-field sector with uniformly high probability.

\subsection{The remaining large-field interface}
\label{fwp-scope}

The distinction between curvature and retained-coordinate windows is
essential. For a link with distinct endpoints, gauge invariance at one
endpoint makes its marginal left-invariant under \(G\). Hence
\begin{equation}
 \mu_g\{U_e\in A\}=\operatorname{Haar}(A)
 \label{fwp-haar-marginal}
\end{equation}
for every measurable \(A\), independently of \(g\).
An unfixed link is therefore outside a shrinking identity window with
probability tending to one. Pure-gauge fields can have all plaquettes
equal to \(I\) and individual links far from \(I\).

The actual retained fluctuation \(\xi\) instead depends on the nonlinear
constraint, its minimizing background, the contour gauge and the pivot
chart. No result above proves that \(\|\xi\|>R\) forces the interior
Wilson action to exceed \(c g^2R^2\). Such a deterministic, quantitative
implication for the exact chart would be an additional input.
A local tree-gauge estimate from small plaquettes does not by itself
identify its potential with \(\xi\), or establish the required
\(C^{1,\alpha}\) background domain.

Global holonomy is a second obstruction. For example,
\(U_1(x)=e^{i\theta H/n_1}\), \(U_\mu(x)=I\) for \(\mu\ne1\),
has every plaquette equal to \(I\), while its noncontractible Wilson
line is \(e^{i\theta H}\). Plaquette rarity cannot simply discard these
sectors.

Finally, the proof uses the positive bare Wilson measure. It has not
been extended here to complex source measures, arbitrary conditioned
boundaries or effective actions after many RG steps. It supplies a
genuine full-measure large-field input: uniform prescribed-event and
cluster bounds on the original regulator. It does not yet supply the
exact retained-window implication, a convergent sum over all large-field
components with analytic source and locality bounds, or a complete
Wilson RG history.

For context, the general reflection/chessboard framework is
Fr\"ohlich--Israel--Lieb--Simon, \emph{Communications in Mathematical
Physics} \textbf{62} (1978), 1--34,
DOI \(10.1007/BF01940327\); gauge positivity is treated by
Osterwalder--Seiler, \emph{Annals of Physics} \textbf{110} (1978),
440--471, DOI \(10.1016/0003\text{-}4916(78)90039\text{-}8\).
The site-plane gluing and normalization used here were proved explicitly
above rather than imported with an unspecified boundary convention.

